\documentclass[11pt]{article}
\usepackage[preprint]{acl}
\usepackage{times}
\usepackage{latexsym}
\usepackage[T1]{fontenc}
\usepackage[utf8]{inputenc}
\usepackage{microtype}
\usepackage{graphicx}
\usepackage{booktabs}
\usepackage{amsmath,amssymb,amsthm}
\usepackage{multirow}
\usepackage{tikz}
\usetikzlibrary{arrows.meta,positioning,fit,backgrounds}
\usepackage{enumitem}

\makeatletter
\ifacl@linenumbers
  \font\aclhvsmall=phvb at 6.5pt
  \renewcommand\linenumberfont{\aclhvsmall\color{lightgray}}
  \setlength{\linenumbersep}{8pt}
\fi
\makeatother

\newtheorem{proposition}{Proposition}
\newcommand{\method}{FinPlanRAG}
\newcommand{\anchor}{\textsc{AnchorEvolve}}
\newcommand{\E}{\mathbb{E}}
\newcommand{\Prob}{\mathbb{P}}
\newcommand{\indic}{\mathbf{1}}

% Generated by make_paper_tables.py
\newcommand{\DSVAcc}{.533}
\newcommand{\DSVBal}{.491}
\newcommand{\DSVMacro}{.486}
\newcommand{\DSVLoss}{3.454}
\newcommand{\DSVBestAcc}{.551}
\newcommand{\DSVBestBal}{.461}
\newcommand{\DSVBestMacro}{.464}
\newcommand{\DSVBestLoss}{4.678}
\newcommand{\DSBestBaseName}{FinDVer}
\newcommand{\DSRouteAcc}{.995}
\newcommand{\DSRouteUtil}{.403}
\newcommand{\DSAlwaysRAG}{.121}
\newcommand{\DSAlwaysAnchor}{.282}
\newcommand{\DSOracle}{.403}
\newcommand{\DSDocAcc}{.457}
\newcommand{\DSDocParas}{10.69}
\newcommand{\DSDocFixedAcc}{.390}
\newcommand{\GLMVAcc}{.533}
\newcommand{\GLMVBal}{.491}
\newcommand{\GLMVMacro}{.486}
\newcommand{\GLMVLoss}{3.454}
\newcommand{\GLMVBestAcc}{.498}
\newcommand{\GLMVBestBal}{.439}
\newcommand{\GLMVBestMacro}{.442}
\newcommand{\GLMVBestLoss}{5.233}
\newcommand{\GLMBestBaseName}{FinDVer}
\newcommand{\GLMRouteAcc}{.990}
\newcommand{\GLMRouteUtil}{.419}
\newcommand{\GLMAlwaysRAG}{.130}
\newcommand{\GLMAlwaysAnchor}{.292}
\newcommand{\GLMOracle}{.422}
\newcommand{\GLMDocAcc}{.497}
\newcommand{\GLMDocParas}{10.31}
\newcommand{\GLMDocFixedAcc}{.480}
\newcommand{\QWVAcc}{.533}
\newcommand{\QWVBal}{.491}
\newcommand{\QWVMacro}{.486}
\newcommand{\QWVLoss}{3.454}
\newcommand{\QWVBestAcc}{.511}
\newcommand{\QWVBestBal}{.477}
\newcommand{\QWVBestMacro}{.464}
\newcommand{\QWVBestLoss}{4.846}
\newcommand{\QWBestBaseName}{Fin-Fact}
\newcommand{\QWRouteAcc}{.867}
\newcommand{\QWRouteUtil}{.263}
\newcommand{\QWAlwaysRAG}{.018}
\newcommand{\QWAlwaysAnchor}{.298}
\newcommand{\QWOracle}{.316}
\newcommand{\QWDocAcc}{.093}
\newcommand{\QWDocParas}{5.85}
\newcommand{\QWDocFixedAcc}{.100}


\title{\method: Routing Mixed Financial Text Queries to Obligation Retrieval and Anchor Evolution}
\author{Anonymous ACL submission}

\begin{document}
\sloppy
\maketitle

\begin{abstract}
Financial text analysis mixes determinate questions with forward-looking claims. The first asks for a fact stated at the cutoff; the second asks about an outcome not yet realized. Existing retrieval-augmented generation and claim verification fix the output contract before seeing the evidence, so they cannot serve both. We present \method, which uses an LLM router to decide whether a query is a forward-looking claim and routes determinate queries to an obligation-retrieval branch and forward-looking claims to an anchor-evolution branch. The determinate branch binds typed obligations, acquires point-in-time evidence until the obligations are covered, and returns a closed answer. The AnchorEvolve branch, \anchor, extracts stable-relation, schema, and instance anchors, qualifies them, updates a belief, and returns a point answer, a credible set, or an abstention under an asymmetric loss.

The paper contributes a routing formulation for mixed financial queries; a financial instantiation of \anchor{} with anchor extraction, qualification, a calibrated posterior, and a loss-aware gate; and a unified evaluation on DocFinQA, XBRL-verified guidance claims, and FinDVer across backbones, named baselines, and metrics. The determinate branch improves accuracy at lower evidence cost; the AnchorEvolve branch preserves accuracy while lowering decision loss and improving calibration; the router approaches oracle mixed utility on the larger backbones. The framework makes the answer form a testable part of the model rather than a fixed interface, and offers an auditable approach to point-in-time financial question answering and claim adjudication.
\end{abstract}

\section{Introduction}

A financial analyst works with two kinds of inputs. A determinate question asks for a fact that a filing states at the cutoff, for example the net revenue change reported in a 10-K or the segment operating profit for a fiscal year. A forward-looking claim asserts something forward-looking, for example that a new production line will start in the second half or that quarterly revenue will land inside management's guidance range. The first admits a closed answer read from evidence. The second does not: at the cutoff the outcome has not happened, and the correct output may be a qualified judgment, a credible set of outcomes, or an abstention.

Existing systems collapse the distinction. Retrieval-augmented generation retrieves passages and generates an answer for every query \citep{lewis2020rag}. Financial claim verifiers return supported, refuted, or insufficient for every claim \citep{zhao2024findver,rangapur2025finfact}. Selective classifiers return an answer or an abstention \citep{chow1970optimum,geifman2017selective}. Adaptive retrieval decides when or how much to retrieve \citep{jiang2023flare,jeong2024adaptiverag}. Each fixes the output contract before seeing the evidence. A verifier applied to a determinate question wastes the chance to return a precise value; a reader applied to a forward-looking claim produces an unqualified point answer with no set and no calibrated credence.

This paper treats the choice of contract as a routing problem. The key insight is that the output contract is itself a decision variable: once the query type is known, each branch chooses the output granularity from the evidence it acquires, rather than fixing that granularity before retrieval. \method{} sends each query to one of two specialist branches. The router is an LLM call that returns a route, a confidence, and a short reason. The obligation-retrieval branch binds a typed obligation tuple (entity, form, fiscal period, metric, unit) to a point-in-time corpus, retrieves evidence under the obligation, and asks an LLM reader for a closed answer. The AnchorEvolve branch, \anchor, extracts anchors at three levels: stable relations (L0), the industry schema (L1), and anchor instances (L2). It qualifies each anchor, updates a belief over outcome buckets, and applies a loss-aware gate that returns a point answer, a credible set, or an abstention. The two branches share a router but not an output contract. Both branches call an LLM, and the router is a separate LLM call.

\paragraph{Contributions.}
\begin{itemize}[leftmargin=*,itemsep=1pt,topsep=2pt]
\item \textbf{Architecture.} We formulate financial question answering as a mixed-workload routing problem over a determinate retrieval contract and a forward-looking claim-adjudication contract. The router, the reader, and the anchor extractor are LLM calls; no branch is a keyword rule.
\item \textbf{Vague branch.} We instantiate \anchor{} for revenue guidance with L0/L1/L2 anchors, LLM extraction, hard/soft/vacuous/adversarial/stale qualification, a belief update, an anchor-conditioned ensemble with a calibrated non-LLM posterior, and a loss-aware gate. The branch returns a point answer, a credible set, or an abstention.
\item \textbf{Determinate branch.} We instantiate the obligation-retrieval branch with typed obligations, point-in-time BM25 retrieval, obligation coverage, and an iterative acquisition planner that stops when the obligations are covered or the reader is confident.
\item \textbf{Evidence.} We evaluate on DocFinQA \citep{reddy2024docfinqa}, 227 XBRL-verified revenue-guidance claims, and FinDVer \citep{zhao2024findver}, with DeepSeek, GLM 5.3 Flash, and a locally served Qwen3-8B. We compare against six source-grounded named baselines. The AnchorEvolve branch improves balanced accuracy and macro F1 in point estimate over every named baseline on the primary backbone, and the router reaches oracle-level mixed utility.
\end{itemize}

\section{Related Work}

\paragraph{Financial question answering and verification.}
Financial QA benchmarks test numerical reasoning over filings, tables, and reports \citep{chen2021finqa,zhu2021tatqa,islam2023financebench,liu2024finben}. DocFinQA adds full 10-K contexts to FinQA questions and tests long-document evidence selection \citep{reddy2024docfinqa}. Claim-verification benchmarks test whether a statement is supported, contradicted, or unverifiable \citep{zhao2024findver,rangapur2025finfact}. Recent financial RAG systems add typed evidence, period constraints, provenance, and corpus-aligned retrieval \citep{mansouri2026finrank,jiang2026finrate,chen2026hcrag,chi2026finsagent}. These systems improve the evidence that reaches a reader, but the output contract is fixed at query time.

\paragraph{Retrieval-augmented generation.}
RAG combines a retriever with a generator and assumes that an answer can be assembled from retrieved context \citep{lewis2020rag}. Active and adaptive variants decide when to retrieve or how much context to retrieve \citep{jiang2023flare,jeong2024adaptiverag}. Agent frameworks interleave reasoning, retrieval, and self-correction \citep{yao2023react,shinn2023reflexion}. These methods adapt the retrieval process, but the output remains a generated answer. They do not represent a credible set, an abstention, or a decision to qualify a claim.

\paragraph{Selective prediction, calibration, and conformal sets.}
The reject option trades recognition error against rejection rate \citep{chow1970optimum}. Selective classification learns a confidence threshold and abstains on low-confidence inputs \citep{geifman2017selective}. Modern classifiers are often miscalibrated, and temperature scaling is a strong baseline \citep{guo2017calibration}. Cost-sensitive learning assigns different prices to different errors \citep{elkan2001cost}. Split conformal prediction converts a score into a set with distribution-free coverage under exchangeability \citep{vovk2005algorithmic,angelopoulos2021gentle}. These components supply the loss, confidence, and set ingredients. They assume a fixed task and a fixed output contract.

\paragraph{Anchors and point-in-time evidence.}
Partial identification represents the outcomes consistent with the available evidence \citep{manski2003partial}. Financial machine learning is vulnerable to lookahead bias when a model uses information that was not available at the decision time \citep{lopezdeprado2018advances}. \anchor{} borrows the idea of constraints and feasible sets but applies it to a claim: an anchor is a relation between a source and a hypothesis, and a hard anchor can rule out a bucket while a soft anchor only reweights. The evidence filter and the schema filter are both point-in-time.

\section{FinPlanRAG Overview}
\label{sec:architecture}

\subsection{Workload and router}
A query $q$ arrives with a cutoff $c$. The workload is a mixture of determinate questions $q\in\mathcal Q_D$ and forward-looking claims $q\in\mathcal Q_V$. A router $r(q,c)\in\{D,V\}$ selects the branch. The router is an LLM call with a short prompt and a JSON response:
\begin{equation}
r(q,c)=\bigl(\text{route},\ \text{confidence},\ \text{reason}\bigr).
\end{equation}
The prompt defines $D$ as a request for a fact, number, passage, estimate, or projection that can be read or computed from a filing at the cutoff, including a question about a future value when the filing states that value. It defines $V$ as a forward-looking claim whose truth depends on an outcome not realized at the cutoff. The route is parsed from JSON; an invalid response falls back to a conservative rule based on the presence of a guidance range and a target period.

The router sees the query text and the cutoff metadata. It does not see the realized outcome or the branch outputs. In deployment it can be cached per query. The cost is one short LLM call, reported separately from the branch cost.

\subsection{Determinate branch: obligation retrieval}
The determinate branch treats a question as a typed obligation. An obligation is a tuple
\begin{equation}
o=(e,\,f,\,y,\,p,\,m,\,u),
\end{equation}
where $e$ is the entity, $f$ the form or document type, $y$ the fiscal year, $p$ the fiscal period, $m$ the metric, and $u$ the unit. The parser derives the obligation terms from the question and the query metadata. A document is admissible only if its availability date is at or before $c$ and its metadata is consistent with the obligation. The branch builds paragraph and table chunks, indexes them with BM25, and runs an iterative acquisition loop.

At round $t$, the branch retrieves a candidate set for the question plus the obligation terms that are not yet covered. It adds up to two new paragraphs, calls the LLM reader on the acquired evidence, and stops when the reader reports sufficient evidence and all obligation terms are covered, or when the reader is confident and no term is missing, or when the budget $B$ is exhausted. The reader returns a JSON object with an answer, a confidence, a sufficiency flag, the missing information, the evidence ids, and a short rationale. The output is the answer when the reader is sufficient; otherwise the branch returns \texttt{insufficient}. The default budget is $B=8$ paragraphs.

\subsection{Vague branch: anchored adjudication}
The AnchorEvolve branch treats a claim as a belief over a finite hypothesis space. In the general \anchor{} design the space is
\begin{equation}
\mathcal H=\{(t,s): t\in\mathcal T,\ s\in\mathcal S\},
\end{equation}
where $t$ is a timing and $s$ a stage. For revenue guidance we instantiate $\mathcal H=\{\text{below},\text{within},\text{above}\}$ relative to the guided range. The branch returns
\begin{equation}
J(q,c)=(h^*,\,b^*,\,R_\alpha,\,\text{action}),
\end{equation}
where $h^*$ is the top bucket, $b^*$ its credence, $R_\alpha$ a $1-\alpha$ credible set, and the action is answer, adjudicate, or abstain. Section~\ref{sec:anchor} gives the construction.

\begin{figure*}[t]
\centering

\resizebox{\textwidth}{!}{%
\begin{tikzpicture}[
  font=\footnotesize,
  >=Latex,
  box/.style={draw, rounded corners=2pt, align=center, text width=2.75cm, minimum height=0.9cm, inner sep=3pt},
  smallbox/.style={draw, rounded corners=2pt, align=center, text width=2.35cm, minimum height=0.72cm, inner sep=2pt},
  arrow/.style={->, thick, draw=black!70}
]
\node[box, fill=blue!6] (q) at (0,0) {Query $q$\\cutoff $c$};
\node[box, fill=yellow!12] (r) at (3.4,0) {LLM router\\route / confidence};
\node[box, fill=green!8] (d) at (7.0,1.6) {Obligation\\Retrieval};
\node[box, fill=orange!10] (v) at (7.0,-1.6) {Forward-looking\\\anchor};
\node[smallbox, fill=green!8] (d1) at (10.8,2.2) {typed obligation\\point-in-time corpus};
\node[smallbox, fill=green!8] (d2) at (10.8,1.1) {BM25 + coverage\\LLM reader};
\node[smallbox, fill=orange!10] (v1) at (10.8,-1.1) {LLM anchor extraction\\L0 / L1 / L2};
\node[smallbox, fill=orange!10] (v2) at (10.8,-2.2) {belief + gate\\answer / set / abstain};
\node[box, fill=blue!6] (out) at (14.4,0) {answer, set,\\or abstention};
\draw[arrow] (q) -- (r);
\draw[arrow] (r) -- node[above, font=\scriptsize]{$D$} (d);
\draw[arrow] (r) -- node[below, font=\scriptsize]{$V$} (v);
\draw[arrow] (d) -- (d1); \draw[arrow] (d) -- (d2);
\draw[arrow] (v) -- (v1); \draw[arrow] (v) -- (v2);
\draw[arrow] (d1) -- (out); \draw[arrow] (d2) -- (out);
\draw[arrow] (v1) -- (out); \draw[arrow] (v2) -- (out);
\end{tikzpicture}
}
\caption{\method{} routes a query to an obligation-retrieval branch or a forward-looking AnchorEvolve branch. The router, the reader, and the anchor extractor are LLM calls. The two branches return different answer forms from the same workload.}
\label{fig:architecture}
\end{figure*}

\section{The \anchor{} Branch}
\label{sec:anchor}

\subsection{Anchor layers}
\anchor{} separates three layers. L0 is a stable relation ontology: bound, precedence, lead time, duration, conservation, rate, unit conversion, and scope. L0 does not change across industries. L1 is an industry schema: each anchor type is a tuple (name, relation in L0, fields, valid-from round). For revenue guidance the initial schema contains guidance width, current-quarter growth, sector history, company history, and other. L2 is an anchor instance: an id, a type, a value or range, a unit, a scope, a source, a time, a reliability, and a validity interval. The separation lets a new industry reuse L0 and change L1 and L2, and it makes every anchor auditable at the source and time level.

\subsection{LLM anchor extraction}
The extractor is an LLM call. It receives the claim, the company and sector, the cutoff, the target period, the guided range, the current-quarter value and reliability, the observed and sector base rates, and the available evidence items. It returns JSON with a list of anchors and a three-way score:
{\scriptsize
\begin{verbatim}
{"anchors":[{"level":"L2","type":"...",
  "relation":"bound|precedence|...",
  "value":"...","scope":"company|...",
  "reliability":0.0,"hard":false,
  "implied_bucket":"below|within|above",
  "rationale":"..."}],
 "bucket_scores":{"below":0.0,
   "within":0.0,"above":0.0},
 "rationale":"..."}
\end{verbatim}}
The prompt states that the guided range defines the three buckets and is not by itself evidence that the outcome will be within. It asks the model to use the current-quarter value relative to the guidance midpoint, the guidance width, and the observed base rates. An anchor is hard only when it is an explicit constraint with reliability at least .65; management guidance is a soft signal by default. The extractor is not asked for the realized outcome and does not receive it.

\subsection{Qualification and belief update}
Each extracted anchor is qualified before it updates the belief:
\begin{itemize}[leftmargin=*,itemsep=1pt,topsep=2pt]
\item \textbf{Hard.} An anchor with an implied bucket and reliability at least .65 intersects the feasible set with the anchor mask. If the intersection is empty, the anchor is treated as a conflict and does not update.
\item \textbf{Soft.} An anchor with an implied bucket contributes a categorical likelihood. The likelihood puts mass $p_e=.7$ on the implied bucket and spreads $1-p_e$ over the other buckets; reliability enters as the exponent.
\item \textbf{Vacuous.} An anchor whose evidence is below the informativeness threshold is skipped.
\item \textbf{Adversarial.} An anchor whose scope does not match the claim, or which contradicts a hard anchor, is skipped.
\item \textbf{Stale.} An anchor whose evidence is superseded or dated after the cutoff is skipped.
\end{itemize}
For a soft anchor with implied bucket $j$ and reliability $\rho$, the likelihood is
\begin{equation}
L_e(k)=
\begin{cases}
p_e, & k=j,\\
(1-p_e)/(|\mathcal H|-1), & k\ne j,
\end{cases}
\end{equation}
Here $p_e=.7$ is fixed by the soft anchor informativeness of .4, and reliability enters through the exponent $\rho$.
The belief starts from a point-in-time prior $b_0$ formed from development and validation outcomes. Hard anchors set the belief to a point mass on the feasible bucket with a small floor; soft anchors update
\begin{equation}
b(k)\propto b(k)\,L_e(k)^{\rho}.
\end{equation}
The extractor's three-way score is a marginal signal. The branch mixes it with the anchor-updated belief at equal weight. The result is the anchor posterior $b_A$.

\subsection{Anchor-conditioned ensemble}
The anchor posterior is fused with a calibrated non-LLM posterior $b_R$ from a random forest on point-in-time structured features: guidance width, current-quarter value relative to the midpoint, the ratio of guidance to current quarter, prior outcome counts and frequencies, previous-guidance growth and width, target quarter, cutoff year, and sector. The random forest is trained on development and validation claims and calibrated by temperature scaling on company-disjoint folds. The fused belief is
\begin{equation}
b_F=w\,b_A+(1-w)\,b_R,
\end{equation}
with $w=.1$ as a candidate. The validation portfolio compares the fused belief with the calibrated RF posterior and selects the RF posterior for the final point answer, ranking, and decision head;

\subsection{Loss-aware gate}
The gate chooses among answer, adjudicate, and abstain under an asymmetric loss. The answer loss charges more for an overclaim than an underclaim:
\begin{equation}
L_{\text{ans}}(a,y)=
\begin{cases}
0 & a=y,\\
\ell_{\text{over}} & a=\text{above}, y=\text{below},\\
\ell_{\text{under}} & a=\text{below}, y=\text{above},\\
\ell_{\text{wrong}} & \text{otherwise}.
\end{cases}
\end{equation}
The set loss charges a miss and a width penalty:
\begin{equation}
\begin{aligned}
L_{\text{set}}(R,y)&=\indic[y\notin R]\,\ell_{\text{miss}}\\
&\quad+\ell_{\text{width}}\max(0,|R|-1),
\end{aligned}
\end{equation}
and abstention costs $\ell_{\text{abstain}}$. The credible set $R_\alpha$ is the smallest set of buckets whose posterior mass reaches $1-\alpha$. The gate evaluates the expected loss of answer, adjudicate, and abstain and selects the minimum. The default parameters are $\ell_{\text{over}}=12$, $\ell_{\text{under}}=8$, $\ell_{\text{wrong}}=10$, $\ell_{\text{miss}}=8$, $\ell_{\text{width}}=2$, $\ell_{\text{abstain}}=5$, and $\alpha=.1$.

The gate is a portfolio over two decision heads. The anchor-fused posterior is optimized for point discrimination; the calibrated RF posterior is optimized for calibrated decisions. The portfolio compares the two heads on the validation split and selects the head with the lower validation loss, with a switching margin of .15. On the revenue-guidance validation split the RF head has loss 3.30 and the fused head has loss 3.96, so the portfolio selects the RF head. The point answer, the credible set, and the action come from the selected decision head. The anchor posterior supplies evidence qualification and the audit trail; it does not replace the calibrated posterior. Test outcomes are not used for fitting, calibration, weight selection, or portfolio selection.

\subsection{Point-in-time self-evolution}
\anchor{} updates its source and anchor reliabilities after outcomes are revealed. For each used evidence item, the branch records whether the evidence was consistent with the realized bucket, updates the source reliability, updates the anchor-type reliability, and updates the classifier weights. Outcomes from round $r$ update the state only after all claims in round $r$ have been adjudicated, and they are used from round $r+1$. A newly discovered anchor type receives a valid-from round of $r+1$, so the claim that motivated its discovery cannot use it. This rule keeps the self-evolution loop point-in-time. The revenue-guidance experiments report the first round; the production-milestone prototype exercises the multi-round loop.

\section{Experimental Design}
\label{sec:experiments}

\subsection{Data}
\textbf{DocFinQA} \citep{reddy2024docfinqa} augments FinQA questions with full 10-K contexts. We use a fixed sample of 300 test questions and the released gold answers. FinQA program constants are used only as a silver evidence proxy, never as reader input.

\textbf{Revenue guidance.} The benchmark contains 1,005 quarterly revenue-guidance claims from U.S. SEC 8-K filings between 2010 and 2025. Each claim contains a guided range and a cutoff. The realized quarter is taken from SEC XBRL, with a reported-filing fallback when XBRL has no standalone quarter. The label is below, within, or above the guided range. The split is temporal: 529 development claims (2010--2018), 203 validation claims (2019--2021), and 273 test claims (2022--2025). The test split has 126 above, 110 within, and 37 below claims. We report the 227 test claims with XBRL-verified or XBRL-overridden labels; 46 text-only claims are excluded from the main results.

\textbf{FinDVer} \citep{zhao2024findver} testmini contains 700 claim-verification examples with financial reports. We use it as an external transfer set.

\subsection{Backbones}
We run three backbones. DeepSeek Flash with a 1M-token context is served through an Anthropic-compatible endpoint. GLM 5.3 Flash is served through an OpenAI-compatible endpoint. Qwen3-8B is served locally by vLLM on one RTX 4080 SUPER with thinking disabled. All methods in a table use the same backbone, temperature zero, and JSON-only prompts. Parse failures are scored as incorrect. On DocFinQA, Qwen3-8B returns verbose answers that fail the exact-match format; every condition is below .11 exact match. We therefore use DeepSeek and GLM for the determinate comparison and report the Qwen DocFinQA block in Appendix~\ref{app:qwen}.

\subsection{Baselines}
The determinate branch is compared with fixed BM25 budgets of 1, 2, 4, 8, and 16 paragraphs, with Self-RAG-style, CRAG-style, FLARE-style, and Adaptive-RAG-style adaptations, and with the obligation retrieval planner. The four named RAG methods keep their native retrieval/reflection strategy and use the same paragraph chunks, reader, and scoring as the obligation retrieval branch; they are strategy adaptations, not unmodified end-to-end systems. All conditions use the same reader and backbone; only the evidence policy changes.

The AnchorEvolve branch is compared with six named adaptations: FinDVer-style \citep{zhao2024findver}, Fin-Fact-style \citep{rangapur2025finfact}, Self-RAG-style \citep{asai2024selfrag}, CRAG-style \citep{yan2024crag}, FLARE-style \citep{jiang2023flare}, and Adaptive-RAG-style \citep{jeong2024adaptiverag}. Each method follows the original paper description and public interface and is adapted to the same backbone through one shared three-way adapter. The adapter maps the native output to below, within, and above, and supplies the same point-in-time evidence block and the same six retrieved earlier claims. When a learned component is unavailable or incompatible, we replace it with an API call that uses the official criteria and record the substitution. The six methods are adaptations, not unmodified end-to-end runs on their original data.



\subsection{Metrics}
\label{sec:metrics}
For DocFinQA we report exact match, token F1, answer coverage, mean acquired paragraphs, mean input tokens, and silver evidence recall. For revenue guidance we report accuracy, balanced accuracy, macro F1, MCC, AUROC, AUPRC, ECE, Brier score, decision loss, credible-set coverage, and mean set size. Balanced accuracy is the mean of per-class recall. Decision loss uses the parameters in Section~\ref{sec:anchor}. Lower is better for loss, ECE, and Brier. Confidence intervals are paired bootstrap intervals over claims. Accuracy differences use a McNemar test on paired predictions. For the router we report routing accuracy, the confusion matrix, and a mixed-workload utility. The utility of a determinate query is its token F1; the utility of a forward-looking query is $1-L/10$, where $L$ is the decision loss. A wrong route receives zero utility. The oracle routes every query correctly; always-RAG sends every query to the determinate branch; always-anchor sends every query to the AnchorEvolve branch.

\section{Results}
\label{sec:results}

\subsection{Determinate retrieval}
Table~\ref{tab:docfinqa} reports DocFinQA. The final branch runs the planner with a budget of 16 paragraphs and escalates to the fixed-16 evidence set when the reader is not sufficient or its confidence is below .8. On DeepSeek it obtains .457 exact match with 10.69 paragraphs; the fixed-16 reader obtains .390 with 16 paragraphs, and the fixed-8 reader .367 with 8. The planner without escalation obtains .460 with 6.85 paragraphs; escalation uses the larger evidence set on 24\% of questions and buys accuracy that the planner alone does not reach. On GLM the branch obtains .497 with 10.31 paragraphs; the fixed-16 reader obtains .480 with 16, and the fixed-8 reader .440 with 8. The planner without escalation obtains .470 with 5.88 paragraphs. The branch is above the fixed-16 reader on both backbones and uses 5.3 and 5.7 fewer paragraphs. The four named RAG adaptations are below the branch on both backbones: on DeepSeek, Self-RAG-style and CRAG-style obtain .373, Adaptive-RAG-style .363, and FLARE-style .330; on GLM, Self-RAG-style and CRAG-style obtain .447, Adaptive-RAG-style and FLARE-style .393. Obligation retrieval with fallback is .084 above the best named baseline on DeepSeek and .050 above it on GLM, with comparable or lower evidence. It also has the highest answer coverage (.770 and .810), so it abstains less often than the fixed budgets.

The determinate branch is therefore an accuracy-cost improvement over the fixed-16 baseline on both usable backbones: higher exact match and lower evidence. The planner alone is cheaper but leaves a small accuracy gap on GLM; the confidence-gated fallback closes it without paying the full 16-paragraph cost on every question. The Qwen3-8B DocFinQA block is reported in Appendix~\ref{app:qwen}; all conditions are below .11 exact match because the reader returns explanatory text, and we do not use that backbone for the determinate comparison.

\begin{table*}[t]
\centering
\scriptsize
\setlength{\tabcolsep}{4pt}
\caption{DocFinQA long-document question answering. All conditions use the same reader and backbone; only the evidence policy changes. Named RAG methods are same-backbone adaptations (suffix -style). ``Paras'' is the mean number of acquired paragraphs and ``Cov.'' the answer coverage.}
\label{tab:docfinqa}
\begin{tabular}{lrrrrr}
\toprule
Condition & Acc. & F1 & Paras & Tokens & Cov. \\
\midrule
\multicolumn{6}{l}{\textit{DeepSeek}} \\
Fixed-4 & .333 & .192 & 4.00 & 459 & .557 \\
Fixed-16 & .390 & .234 & 16.00 & 1663 & .640 \\
\textit{Obligation + fallback} & .457 & .276 & 10.69 & 1604 & .770 \\
Self-RAG-style & .373 & .223 & 8.03 & 196 & .613 \\
CRAG-style & .373 & .222 & 8.05 & 193 & .613 \\
FLARE-style & .330 & .197 & 4.00 & 195 & .557 \\
Adaptive-RAG-style & .363 & .211 & 7.22 & 192 & .570 \\
\midrule
\multicolumn{6}{l}{\textit{GLM 5.3}} \\
Fixed-4 & .400 & .199 & 4.00 & 1010 & .653 \\
Fixed-16 & .480 & .249 & 16.00 & 3595 & .757 \\
\textit{Obligation + fallback} & .497 & .252 & 10.31 & 2476 & .810 \\
Self-RAG-style & .447 & .231 & 8.16 & 1909 & .717 \\
CRAG-style & .447 & .235 & 8.32 & 1937 & .727 \\
FLARE-style & .393 & .193 & 4.19 & 1048 & .643 \\
Adaptive-RAG-style & .393 & .198 & 6.51 & 1526 & .657 \\
\bottomrule
\end{tabular}
\end{table*}


\subsection{Forward-looking adjudication}
Table~\ref{tab:vague} reports the 227 verified revenue-guidance claims. The validation portfolio compares the calibrated RF head with the fused head and selects the RF head for the point and decision outputs. The anchor extractor runs on all three backbones; the reported FinPlanRAG metrics therefore use a backbone-independent belief head, while the LLM anchor posterior varies by backbone. On DeepSeek, \method{} obtains accuracy .533, balanced accuracy .491, macro F1 .486, and decision loss 3.454. FinDVer-style obtains .551, .461, .464, and 4.678. Plain accuracy is .018 lower, which is a statistical tie. Balanced accuracy is .030 higher and macro F1 is .022 higher; both intervals include zero. Decision loss falls by about 1.2 points, and the paired interval excludes zero.

On GLM, \method{} obtains .533, .491, .486, and 3.454, against .498, .439, .442, and 5.233 for FinDVer-style. Accuracy is .035 higher, balanced accuracy .052 higher, macro F1 .044 higher, and decision loss about 1.8 points lower. On Qwen3, \method{} obtains .533, .491, .486, and 3.454, against .511, .477, .464, and 4.846 for Fin-Fact-style. Accuracy is .022 higher, balanced accuracy .014 higher, macro F1 .022 higher, and decision loss about 1.4 points lower.

The robust gain is decision loss: the paired interval excludes zero on every backbone. Plain accuracy is higher on GLM and Qwen3 and statistically tied on DeepSeek. Balanced accuracy and macro F1 are higher in point estimate on every backbone, but their intervals include zero. The result is a decision-theoretic gain rather than a uniform classification gain: plain accuracy is comparable, and the system prices overclaims, set width, and abstention better than the named baselines.

\begin{table*}[t]
\centering
\scriptsize
\setlength{\tabcolsep}{2.6pt}
\caption{Revenue-guidance adjudication on 227 XBRL-verified test claims. Acc: accuracy; Bal: balanced accuracy; Macro: macro F1; Loss: asymmetric decision loss (lower is better). FinPlanRAG is the full system. Named methods are same-backbone adaptations (suffix -style).}
\label{tab:vague}
\resizebox{0.92\textwidth}{!}{%
\begin{tabular}{lrrrrrrrrrrrr}
\toprule
Method & DeepSeek Acc & Bal & Macro & Loss & GLM 5.3 Acc & Bal & Macro & Loss & Qwen3-8B Acc & Bal & Macro & Loss \\
\midrule
FinDVer-style & \textbf{.551} & .461 & .464 & 4.678 & .498 & .439 & .442 & 5.233 & .445 & .422 & .406 & 5.524 \\
Fin-Fact-style & .537 & .446 & .450 & 4.705 & .454 & .398 & .402 & 5.339 & .511 & .477 & .464 & 4.846 \\
Self-RAG-style & .546 & .451 & .456 & 4.564 & .454 & .406 & .404 & 5.163 & .489 & .447 & .437 & 5.093 \\
CRAG-style & .529 & .423 & .418 & 4.802 & .515 & .431 & .435 & 4.846 & .454 & .418 & .399 & 5.366 \\
FLARE-style & .511 & .417 & .411 & 4.881 & .471 & .388 & .387 & 5.286 & .326 & .298 & .212 & 6.740 \\
Adaptive-RAG-style & .493 & .406 & .400 & 5.093 & .427 & .372 & .376 & 5.630 & .344 & .324 & .290 & 6.485 \\
\midrule
FinPlanRAG & .533 & \textbf{.491} & \textbf{.486} & \textbf{3.454} & \textbf{.533} & \textbf{.491} & \textbf{.486} & \textbf{3.454} & \textbf{.533} & \textbf{.491} & \textbf{.486} & \textbf{3.454} \\
\bottomrule
\end{tabular}
}
\end{table*}


\subsection{Routing and mixed workload}
Table~\ref{tab:router} reports the mixed workload. The DeepSeek router routes 570 of 573 queries correctly (.995): 297 of 300 determinate questions and all 273 forward-looking claims. Three determinate questions are sent to the anchor branch; their F1 is zero, so the router still reaches .403 mixed utility, equal to the oracle, against .121 for always-RAG and .282 for always-anchor. The GLM router routes 567 of 573 correctly (.990): 297 of 300 determinate questions and 270 of 273 forward-looking claims. It reaches .419 mixed utility against .422 for the oracle, .130 for always-RAG, and .292 for always-anchor.

The Qwen3 router is weaker. It routes 271 of 300 determinate questions and 226 of 273 forward-looking claims correctly (.867), and reaches .263 mixed utility. The always-anchor policy obtains .298, so the Qwen router is below the better fixed branch. The errors are balanced in direction (29 determinate questions sent to anchor, 47 claims sent to RAG), which suggests that an 8B backbone does not reliably separate the two contracts. The router result is therefore reported per backbone rather than pooled.

\begin{table}[t]
\centering
\small
\setlength{\tabcolsep}{3pt}
\caption{Mixed-workload routing over 300 DocFinQA questions and 273 forward-looking claims. Utility is the task-specific score defined in Section~\ref{sec:metrics}.}
\label{tab:router}
\resizebox{\columnwidth}{!}{%
\begin{tabular}{lrrrrr}
\toprule
Backbone & Route acc. & Router & Always RAG & Always anchor & Oracle \\
\midrule
DeepSeek & .995 & .403 & .121 & .282 & .403 \\
GLM 5.3 & .990 & .419 & .130 & .292 & .422 \\
Qwen3-8B & .867 & .263 & .018 & .298 & .316 \\
\bottomrule
\end{tabular}
}
\end{table}


\subsection{External transfer}
Table~\ref{tab:findver} reports FinDVer testmini with the default reader. On DeepSeek the planner with fallback obtains .701 accuracy with 8.76 paragraphs; the fixed-16 reader obtains .660 with 16, and the fixed-8 reader .623 with 8. On GLM the branch obtains .753 with 9.19 paragraphs; the fixed-16 reader obtains .743 with 16, and the fixed-8 reader .680 with 8. On Qwen3 the branch obtains .649 with 6.13 paragraphs; the fixed-16 reader obtains .636 with 16. The branch is above the fixed-16 reader on all three backbones and uses 6.9 to 9.9 fewer paragraphs. The official chain-of-thought reader with gold evidence obtains .817 on DeepSeek and .819 on GLM, and with ten retrieved paragraphs .787 on both; these use a different reader and evidence interface and are an upper reference, not a like-for-like baseline. The transfer result supports the same conclusion as DocFinQA: the obligation planner plus confidence-gated escalation is a cost-effective determinate policy, not a uniform accuracy winner over a stronger reader.

\begin{table*}[t]
\centering
\scriptsize
\setlength{\tabcolsep}{4pt}
\caption{FinDVer testmini transfer with the default reader.}
\label{tab:findver}
\begin{tabular}{llrrrr}
\toprule
Backbone & Condition & Acc. & Macro F1 & Paras & Coverage \\
\midrule
\multicolumn{6}{l}{\textit{DeepSeek}} \\
Fixed-4 & .551 & .665 & 4.00 & .659 \\
Fixed-8 & .623 & .713 & 8.00 & .747 \\
Fixed-16 & .660 & .735 & 16.00 & .796 \\
\textit{Obligation planner} & .680 & .752 & 5.63 & .807 \\
\textit{Obligation + fallback} & .701 & .759 & 8.76 & .847 \\
\midrule
\multicolumn{6}{l}{\textit{GLM 5.3}} \\
Fixed-4 & .599 & .701 & 4.00 & .696 \\
Fixed-8 & .680 & .753 & 8.00 & .803 \\
Fixed-16 & .743 & .792 & 16.00 & .874 \\
\textit{Obligation planner} & .726 & .774 & 5.69 & .871 \\
\textit{Obligation + fallback} & .753 & .789 & 9.19 & .906 \\
\midrule
\multicolumn{6}{l}{\textit{Qwen3-8B}} \\
Fixed-4 & .593 & .635 & 4.00 & .864 \\
Fixed-8 & .629 & .664 & 8.00 & .891 \\
Fixed-16 & .636 & .673 & 16.00 & .887 \\
\textit{Obligation planner} & .637 & .657 & 5.12 & .937 \\
\textit{Obligation + fallback} & .649 & .663 & 6.13 & .954 \\
\bottomrule
\end{tabular}
\end{table*}


\section{Discussion}

The experiments support three claims. First, the output contract should depend on the query. A determinate question benefits from typed obligations and a closed reader; a forward-looking claim benefits from anchors, a belief, and a gate. Second, the AnchorEvolve branch needs both a semantic extractor and a calibrated posterior. The LLM anchors alone are miscalibrated, and the calibrated RF posterior is the strongest point head; a small anchor weight does not improve the point metrics. The final system keeps the RF head for point and decision outputs and uses the anchors for evidence qualification and audit. Third, the router is cheap relative to the branches and recovers almost all of the oracle utility on DeepSeek and GLM, while an 8B backbone is not a reliable query-type detector.

The design has limits. The revenue-guidance instantiation is a three-bucket problem; the general twelve-state timing-by-stage space is exercised only in the production-milestone prototype. The anchor extractor can be wrong, and a wrong hard anchor can eliminate the correct bucket. The qualification rules mitigate this with a reliability floor and a conflict check, but they do not make the extractor infallible. The router is evaluated on a constructed mixed workload with clear provenance; ambiguous queries that sit between the two contracts are not covered. The FinDVer transfer uses a different reader interface, so the comparison there is a frontier check rather than a controlled ablation.

The results also clarify the relation to the earlier decision-interface work. The v2 portfolio over fixed adjudication, abstention, and a planner is a special case of the AnchorEvolve branch's gate: it fixes the evidence model and the output contract and selects a policy on validation. \method{} keeps the portfolio as the decision head but adds the anchor extractor, the LLM router, and the determinate branch. The portfolio is therefore a component, not the whole system.

\section{Conclusion}

\method{} routes financial queries to a determinate retrieval branch or a forward-looking AnchorEvolve branch. The determinate branch binds typed obligations and stops when they are covered. The AnchorEvolve branch extracts L0/L1/L2 anchors, qualifies them, updates a belief, and returns an answer, a credible set, or an abstention under an asymmetric loss. On DocFinQA the planner matches the large-budget reader at a fraction of the evidence. On verified revenue guidance the AnchorEvolve branch improves balanced accuracy and macro F1 in point estimate over six named baselines and lowers decision loss. On a mixed workload the router reaches oracle-level utility. The architecture gives financial RAG a way to choose its output contract instead of fixing it in advance.

\bibliography{references}

\appendix

\section{Prompts and Implementation Details}
The router prompt defines the two contracts, gives four examples, and asks for JSON with route, confidence, and reason. The anchor prompt receives the claim, company, sector, cutoff, target period, guided range, current-quarter value and reliability, observed and sector base rates, and evidence items. It asks for L2 anchors with type, relation, value, scope, reliability, hard flag, implied bucket, and rationale, plus a three-way score. The reader prompt asks for a JSON object with answer, confidence, sufficiency, missing information, evidence ids, and a short rationale. The full prompts and the generation scripts are released with the code.

The DocFinQA runner builds paragraph and table chunks, indexes them with BM25, and uses the same reader for every condition. The AnchorEvolve runner stores one JSONL row per claim with the belief, the credible set, the anchors, the action, and the token usage. The router runner stores one JSONL row per query with the true route, the predicted route, the confidence, the reason, and the branch score. All API responses are cached by a hash of the request payload.

\section{Additional Metrics}
Table~\ref{tab:fullmetrics} reports the full metric suite for the primary backbone. \method{} has a higher AUROC, AUPRC, and balanced accuracy than every named baseline, the lowest ECE, and MCC within .004 of the best baseline. The validation portfolio selects the calibrated RF head for the final outputs. FinDVer-style has higher plain accuracy by .018, but \method{} has a lower Brier score by .014; the accuracy difference is not significant. The reported decision loss is 3.454.

\begin{table*}[t]
\centering
\scriptsize
\setlength{\tabcolsep}{4pt}
\caption{Full metric suite on the 227 verified revenue-guidance claims, DeepSeek backbone. Lower is better for ECE, Brier, decision loss, and AURC; Sel@50 is selective accuracy at 50\% coverage.}
\label{tab:fullmetrics}
\begin{tabular}{lrrrrrrrrrrr}
\toprule
Method & Acc. & Bal. & Macro & MCC & AUROC & AUPRC & ECE & Brier & Loss & Sel@50 & AURC \\
\midrule
FinDVer-style & .551 & .461 & .464 & .233 & .674 & .494 & .092 & .580 & 4.678 & .675 & .337 \\
Fin-Fact-style & .537 & .446 & .450 & .201 & .650 & .475 & .075 & .590 & 4.705 & .675 & .338 \\
Self-RAG-style & .546 & .451 & .456 & .216 & .670 & .491 & .102 & .580 & 4.564 & .684 & .333 \\
CRAG-style & .529 & .423 & .418 & .177 & .643 & .473 & .097 & .593 & 4.802 & .658 & .349 \\
FLARE-style & .511 & .417 & .411 & .164 & .634 & .467 & .130 & .610 & 4.881 & .596 & .362 \\
Adaptive-RAG-style & .493 & .406 & .400 & .143 & .584 & .425 & .147 & .634 & 5.093 & .579 & .415 \\
\method & .533 & .491 & .486 & .229 & .706 & .505 & .033 & .566 & 3.454 & .649 & .383 \\
\bottomrule
\end{tabular}
\end{table*}

\section{Reproducibility}
The released artifacts contain the per-claim predictions, the per-query router decisions, the DocFinQA per-example rows, the analysis scripts, and the table generator. The benchmark data are derived from public SEC filings and XBRL. The named baselines are run through the shared adapter described in Section~\ref{sec:experiments}; the provenance file lists every reused and substituted component. The API keys are read from environment variables and are not stored in the artifacts.


\section{Qwen3-8B on DocFinQA}
\label{app:qwen}
Table~\ref{tab:qwendoc} reports the Qwen3-8B DocFinQA block. The reader returns syntactically valid JSON but often with explanatory text in the answer field, so exact match and token F1 are near floor. The planner still has the highest answer coverage, but the backbone is not informative for the determinate accuracy comparison.

\begin{table}[t]
\centering
\small
\setlength{\tabcolsep}{3pt}
\caption{Qwen3-8B DocFinQA results on the 300-question sample.}
\label{tab:qwendoc}
\begin{tabular}{lrrrr}
\toprule
Condition & Acc. & F1 & Paras & Cov. \\
\midrule
Fixed-1 & .063 & .029 & 1.00 & .590 \\
Fixed-4 & .100 & .048 & 4.00 & .760 \\
Fixed-8 & .110 & .040 & 8.00 & .833 \\
Fixed-16 & .100 & .046 & 16.00 & .853 \\
Obligation planner & .093 & .034 & 3.67 & .883 \\
\bottomrule
\end{tabular}
\end{table}

\end{document}
